% ======================================================================
% 第 VI 节：相变
% ======================================================================
\section{相变}\label{sec:phases}

在强耦合极限（$g\to\infty$，$\kappa\to0$）下，规范理论远非洛伦兹不变的。
更准确地说，由于作用量定义在欧几里得度规上，缺失的是欧几里得不变性。在强耦合
极限下，真空期望值在仅有几个格点间距的分离处就迅速衰减（每分离一个格点间距就
有一个因子 $g^{-2}$ 或 $\kappa$ 或两者兼有）。这对应于存在远大于截断的质量。
[通常的规则是：如果传播子对大的 $x$ 按 $e^{-\xi x}$ 衰减，那么对传播子有贡献的
最低质量中间态的质量为 $1/\xi$。如果传播子对距离 $x=na$ 按 $g^{-2x/a}$ 行为，
那么相应的质量为 $2(\ln g)/a$。若 $g$ 很大，这大于截断动量 $\pi/a$。]

因此，实践中感兴趣的是使关联长度 $\xi$ 远大于格点间距 $a$ 的 $g$ 和 $\kappa$ 值，
以便相应的质量远小于截断。从统计力学知道，大的关联长度与二阶相变（临界点）相联
系。因此，人们寻求特殊的 $g$ 和 $\kappa$ 值 $g_{c}$ 与 $\kappa_{c}$，使 $g$ 和
$\kappa$ 处发生相变。

前面已经论证，规范场存在两个截然不同的相：大 $g$ 的强耦合相束缚夸克，小 $g$ 的
弱耦合相不束缚夸克。这些论证忽略了夸克真空圈，这在 $\kappa$ 小时是合理的。这两
个相之间应当存在转变，它发生在临界值 $g_{c}$ 处，对任何 $g$ 和任何小的 $\kappa$
值都是如此。这是一个可能的相变；正是这一相变在第二节讨论过。但是，正如下面将论证
的，这很可能是一阶相变而非二阶相变。

假设希望用本文的格点理论构造一个强相互作用模型，其中规范群与普通的
SU(3)$\times$SU(3) 对称性分离。那么规范场将全是 SU(3)$\times$SU(3) 单态，而夸克
场除携带规范群量子数外还携带 SU(3)$\times$SU(3) 量子数。稍加思考即可看出，在
强耦合极限（$g\to\infty$，$\kappa\to0$）下，SU(3)$\times$SU(3) 是精确对称性而非
自发破缺的对称性。改变 $g$ 不会改变这一状况；因此必须希望，通过增大 $\kappa$ 能
把精确的 SU(3)$\times$SU(3) 变成破缺的 SU(3)$\times$SU(3)。如果这行不通，可以自由
地在夸克场作用量中加入额外项，以期迫使 SU(3)$\times$SU(3) 自发破缺。为简单起见，
假定 SU(3)$\times$SU(3) 可以通过增大 $\kappa$ 而破缺。那么对从精确的
SU(3)$\times$SU(3) 变为自发破缺的 SU(3)$\times$SU(3) 的 $\kappa$，在临界值
$\kappa_{c}$ 处将发生相变。如果这一相变是二阶的，那么 $\kappa$ 在 $\kappa_{c}$
附近将有大的关联长度；此时该理论可能成为破缺 SU(3)$\times$SU(3) 的一个现实模型，
条件是 $\kappa$ 略大于 $\kappa_{c}$（且 $g$ 足够大以维持夸克束缚）。

总之，真正令人感兴趣的相变是 $\kappa$（或引入夸克作用量的某个其它参数）中的相变，
而不是 $g$ 中的相变。

除特殊极限（$g\to\infty$ 和 $\kappa\to0$，或 $g$ 很小）外，求解格点理论非常困难。
尤其是在关联长度很大的临界点附近求解格点理论更加困难。统计力学工作者已经发展出
各种方法处理这一问题。本节余下部分将简要讨论这些方法。基本有三种途径：
(1) 平均场技术，(2) 级数展开，和 (3) 重正化群方法。

平均场技术${}^{15}$是研究临界点最简单、最粗略的方法；在研究生疏情形时人们总是
首先使用它们。它们用于确定是否存在相变、相变是一阶还是二阶，并粗略估计相变附近
的行为。平均场计算的任何结果都不是完全可信的。稍后将给出平均场计算的例子。

级数展开的一个例子是：对小的动量（动量 $\ll 1/a$），把流-流传播子按 $g$ 和
$\kappa$ 展开，展开到 $g^{-2}$ 和 $\kappa$ 的高阶。然后用 Pad\'e 近似技术寻找
$g$ 或 $\kappa$ 中与质量趋于 0 相联系的奇点。在简单的统计力学问题中，类似的展开
可以生成 12 项左右。本文格点理论的展开更复杂，但经过练习或可望生成约 6 或 7 阶。
级数展开比平均场计算费力得多；它们主要适用于传播子，对三点和四点函数进行级数
展开非常笨拙，而且在尝试此类计算之前必须清楚自己试图学到什么。最佳级数展开
形式之一见文献 16；一般综述见文献 13。

重正化群方法有潜力成为研究临界点附近格点理论最强大、最精确的方法，但目前重正化
群技术的适用范围太有限，尚不能应用于当前问题。见文献 6 和 17。

回到平均场思想。典型的平均场计算是对伊辛铁磁体计算作为外场函数的磁化强度。
令 $s_{n}$ 为格点 $n$ 处的自旋，只取值 $\pm 1$；令相互作用为
\begin{equation}
A = K\sum_{n,\mu}s_{n}s_{n+\hat\mu} + h\sum_{n}s_{n},
\label{eq:6.1}
\end{equation}
其中 $K$ 与自旋-自旋耦合有关，$h$ 正比于外场。于是
\begin{equation}
M = Z^{-1}\Bigl(s_{0}\exp\Bigl[K\sum_{n,\mu}s_{n}s_{n+\hat\mu}
    + h\sum_{n}s_{n}\Bigr]\Bigr),
\label{eq:6.2}
\end{equation}
其中 $(\ )$ 表示对所有自旋的所有组态求和，而
\begin{equation}
Z = \sum\exp\Bigl[K\sum_{n,\mu}s_{n}s_{n+\hat\mu}
    + h\sum_{n}s_{n}\Bigr].
\label{eq:6.3}
\end{equation}

在平均场近似中，假定与 $s_{0}$ 耦合的自旋 $s_{n,\hat\mu}$ 可以用它们的平均值 $M$
代替。结果，$M$ 的公式化简为只对 $s_{0}$ 求和，即
\begin{equation}
M = Z_{0}^{-1}\sum_{s_{0}=\pm 1}s_{0}\exp\bigl[(2dKM+h)s_{0}\bigr],
\label{eq:6.4}
\end{equation}
其中 $d$ 是维数（通常为 3），而
\begin{equation}
Z_{0} = \sum_{s_{0}=\pm 1}\exp\bigl[(2dKM+h)s_{0}\bigr].
\label{eq:6.5}
\end{equation}
结果是
\begin{equation}
M = \tanh(2dKM+h).
\label{eq:6.6}
\end{equation}
如果 $2dK<1$，该方程对作为 $h$ 函数的 $M$ 有唯一解；特别是 $h=0$ 时 $M=0$。
如果 $2dK>1$，解是多值的；稳定性考虑表明，$h=0$ 时必须选取 $M\ne0$ 的解。

在这一近似中，实际上已把 $\sum_{\mu}s_{n,\hat\mu}$ 换成 $2dM$，如果 $d$ 很大这是
好的近似。这通常是平均场理论的性质。

可以对格点规范理论进行类似的平均场计算。此时简单的外场项具有形式
\[
\frac{1}{g}\sum_{n,\mu}\bigl(e^{i\thl{0\mu}(n)}+e^{-i\thl{0\mu}(n)}\bigr)
\]
（这一项要加到规范场作用量上），而 $M$ 可以定义为 $e^{i\thl{0\mu}}$ 的期望值。
问题在于 $M$ 在极限 $h\to0$ 下是否为零。在规范场理论中，$e^{i\thl{0\mu}}$ 与
另外三个指数的乘积耦合；作为平均场近似，把这一乘积用 $M^{3}$ 代替。结果是
\begin{equation}
M = Z_{\theta}^{-1}\int\ddth{0\mu}\,e^{i\thl{0\mu}}
  \exp\bigl[2(d-1)M^{3}+h\bigl(e^{i\thl{0\mu}}+e^{-i\thl{0\mu}}\bigr)\bigr],
\label{eq:6.7}
\end{equation}
其中 $d$ 是时空维数，
\begin{equation}
Z_{\theta} = \int\ddth{0\mu}
  \exp\bigl[2(d-1)M^{3}+h\bigl(e^{i\thl{0\mu}}+e^{-i\thl{0\mu}}\bigr)\bigr],
\label{eq:6.8}
\end{equation}
结果是
\begin{equation}
M = f\bigl[2(d-1)M^{3}+h\bigr],
\label{eq:6.9}
\end{equation}
其中 $f$ 是贝塞尔函数的比值。如果 $g$ 很大，该方程的解唯一且 $h=0$ 时 $M=0$。
如果 $g$ 很小，那么 $h=0$ 时存在 $M\ne0$ 的解，稳定性考虑再次表明应优先选取
$M\ne0$ 的解。

在磁学情形中，发现自发磁化强度 $M$ 在 $2dK-1\to0$ 时趋于零 [见
式~\eqref{eq:6.6}]。然而，规范场情形在 $h=0$ 时从来没有 $M$ 小而非零的解。因此
在 $M$ 由零变为非零的 $g$ 值处存在一阶相变。

在极限 $h\to0$ 下 $M$ 非零意味着规范场对称性发生自发破缺。所以对小 $g$，理论
显示出自发破缺。

关于平均场近似更彻底的讨论由 Balian、Drouffe 和 Itzykson 给出。${}^{17}$ 格点
规范理论的哈密顿量形式由 Kogut 和 Susskind 给出。${}^{18}$ 格点理论中夸克禁闭
的清晰综述见文献 20。强耦合规范理论与弦模型之间联系的另一种表述见文献 21。
